\documentclass{article}

\usepackage{amsmath,amssymb}
\usepackage{xurl}
\usepackage[hidelinks]{hyperref}
\setlength{\emergencystretch}{2em}

\input{command}

\title{Nil-Matrix Product-Length Bound in the MPU Blocking Argument}
\author{}
\date{2026-08-08}

\begin{document}
\maketitle
\gapnote{local-correction}{resolved}

\section{The source assertion}

Let $A$ be a possibly nonunital associative subalgebra of the endomorphisms of
a complex vector space of dimension $D^2$, and suppose that every element of
$A$ is nilpotent. The matrix-product-unitary blocking argument invokes the
Nagata--Higman and Razmyslov results to choose a mixed-product length $J'$
such that
\begin{align}
    J'
        &<D^2,
    \label{eq:mpu_nil_matrix_bound_source_strict_length}\\
    A^{J'}
        &=0.
    \label{eq:mpu_nil_matrix_bound_source_product_vanishing}
\end{align}
Here $A^{J'}=0$ means that every ordered product of $J'$ elements of $A$
vanishes. This assertion appears in lines~405--426 of
Ref.~\cite{Cirac2017MPU}.

The references cited there are Nagata, \emph{Journal of the Mathematical
Society of Japan} 4 (1952), 296--301; Higman, \emph{Proceedings of the
Cambridge Philosophical Society} 52 (1956), 1--4; and Razmyslov,
\emph{Izvestiya Akademii Nauk SSSR, Seriya Matematicheskaya} 38 (1974).
Their nilpotency theorems and general quantitative bounds do not supply the
strict linear estimate
(\ref{eq:mpu_nil_matrix_bound_source_strict_length}).

\section{The corrected nil-matrix bound}

Let $V$ be an $n$-dimensional vector space over a field, and let
$A\subseteq\operatorname{End}(V)$ be a possibly nonunital associative
subalgebra whose elements are all nilpotent. The finite-dimensional nil-matrix
theorem, also called Levitzki's theorem and obtainable from Engel's theorem by
simultaneous strict triangularization, gives
\begin{align}
    A^n
        &=0.
    \label{eq:mpu_nil_matrix_bound_dimension_product_vanishing}
\end{align}
Applying
(\ref{eq:mpu_nil_matrix_bound_dimension_product_vanishing}) to the
$D^2$-dimensional transfer space gives
\begin{align}
    J'
        &\leq D^2.
    \label{eq:mpu_nil_matrix_bound_corrected_transfer_length}
\end{align}

The weak inequality is sharp in general. The algebra of strictly
upper-triangular $n\times n$ matrices contains a nonzero ordered product of
$n-1$ matrix units. Therefore no general argument can replace
(\ref{eq:mpu_nil_matrix_bound_corrected_transfer_length}) by the strict bound
in (\ref{eq:mpu_nil_matrix_bound_source_strict_length}).

\section{Effect on the blocking estimate}

The source has already obtained a first Jordan-block length $J$ satisfying
\begin{align}
    J
        &<D^2.
    \label{eq:mpu_nil_matrix_bound_first_jordan_length}
\end{align}
Combining
(\ref{eq:mpu_nil_matrix_bound_first_jordan_length}) with
(\ref{eq:mpu_nil_matrix_bound_corrected_transfer_length}) yields
\begin{align}
    JJ'
        &<D^2D^2,
    \label{eq:mpu_nil_matrix_bound_product_estimate}\\
    JJ'
        &<D^4.
    \label{eq:mpu_nil_matrix_bound_final_blocking_estimate}
\end{align}
The final strict blocking estimate is unchanged. Only the intermediate claim
$J'<D^2$ must be replaced by $J'\leq D^2$.

The formal development proves
(\ref{eq:mpu_nil_matrix_bound_dimension_product_vanishing}) for arbitrary
finite-dimensional vector spaces and transports it to square
matrices.\footnote{Formal declarations:
\leanid{NonUnitalSubalgebra.list_prod_eq_zero_of_forall_isNilpotent} and
\leanid{NonUnitalSubalgebra.matrix_list_prod_eq_zero_of_forall_isNilpotent}
in \doclink{TNLean/Algebra/NilMatrixSubalgebra}.}
It includes the zero-dimensional case and fixes the ordered-product
convention. The strict-upper-triangular example explains why a strict bound is
unavailable; no formalization of that example is claimed.

\section{Classification}

\emph{Local correction.} Replacing the unsupported intermediate strict bound
with the finite-dimensional nil-matrix bound preserves every hypothesis and
conclusion required by the surrounding blocking argument. In particular,
(\ref{eq:mpu_nil_matrix_bound_final_blocking_estimate}) retains the source's
final strict estimate.

\bibliographystyle{alpha}
\bibliography{references}

\end{document}
